\documentclass[11pt]{article}
\newcommand{\videoid}{N-8tx7lDRBY}
% Preámbulo compartido de las notas del curso de topología algebraica.
% Cada nota empieza con \documentclass[11pt]{article}, define \videoid{...}
% (el identificador del video de YouTube) y luego hace % Preámbulo compartido de las notas del curso de topología algebraica.
% Cada nota empieza con \documentclass[11pt]{article}, define \videoid{...}
% (el identificador del video de YouTube) y luego hace % Preámbulo compartido de las notas del curso de topología algebraica.
% Cada nota empieza con \documentclass[11pt]{article}, define \videoid{...}
% (el identificador del video de YouTube) y luego hace \input{preambulo}.
% Compilar con LuaLaTeX:  lualatex nombre-de-la-nota.tex
\usepackage[spanish,provide=*]{babel}
\usepackage[T1]{fontenc}
\usepackage{lmodern}
\usepackage[a4paper,margin=2.6cm]{geometry}
\usepackage{amsmath,amssymb,amsthm}
\usepackage{xcolor}
\usepackage{tikz}
\usetikzlibrary{arrows.meta,babel}
\usepackage[colorlinks=true,urlcolor=blue!50!black,linkcolor=blue!50!black]{hyperref}

\hyphenpenalty=10000 \exhyphenpenalty=10000 \emergencystretch=3em

% Colores de las figuras (los mismos que usa la página)
\definecolor{extremo}{RGB}{47,107,79}
\definecolor{azul}{RGB}{44,74,154}

% Definiciones numeradas por unidad
\theoremstyle{definition}
\newtheorem{defn}{Definición}
% Numeración por unidad: \unidad{u}{n} (después de \input{preambulo}) hace que las
% definiciones de la nota se llamen u.n, u.(n+1), ...; n es el primer número
% libre de esa unidad (ver la tabla de CONVENCIONES.md).
\newcommand{\unidad}[2]{\renewcommand{\thedefn}{#1.\arabic{defn}}\setcounter{defn}{\numexpr#2-1\relax}}

% Observaciones y erratas del video. Usar SOLO si hay algo que señalar:
%   \begin{erratas}
%     \errata{203}{3:23}{Se dice ..., pero debe decir ...}
%   \end{erratas}
% El primer argumento son los segundos (enlaza al minuto exacto del video),
% el segundo es la marca visible, el tercero es el texto.
\newenvironment{erratas}{\section*{Observaciones y erratas del video}\begin{itemize}}{\end{itemize}}
\newcommand{\errata}[3]{%
  \edef\errurl{https://www.youtube.com/watch?v=\videoid\string&t=#1}%
  \item[\textbf{\expandafter\href\expandafter{\errurl}{#2}}] #3}
.
% Compilar con LuaLaTeX:  lualatex nombre-de-la-nota.tex
\usepackage[spanish,provide=*]{babel}
\usepackage[T1]{fontenc}
\usepackage{lmodern}
\usepackage[a4paper,margin=2.6cm]{geometry}
\usepackage{amsmath,amssymb,amsthm}
\usepackage{xcolor}
\usepackage{tikz}
\usetikzlibrary{arrows.meta,babel}
\usepackage[colorlinks=true,urlcolor=blue!50!black,linkcolor=blue!50!black]{hyperref}

\hyphenpenalty=10000 \exhyphenpenalty=10000 \emergencystretch=3em

% Colores de las figuras (los mismos que usa la página)
\definecolor{extremo}{RGB}{47,107,79}
\definecolor{azul}{RGB}{44,74,154}

% Definiciones numeradas por unidad
\theoremstyle{definition}
\newtheorem{defn}{Definición}
% Numeración por unidad: \unidad{u}{n} (después de % Preámbulo compartido de las notas del curso de topología algebraica.
% Cada nota empieza con \documentclass[11pt]{article}, define \videoid{...}
% (el identificador del video de YouTube) y luego hace \input{preambulo}.
% Compilar con LuaLaTeX:  lualatex nombre-de-la-nota.tex
\usepackage[spanish,provide=*]{babel}
\usepackage[T1]{fontenc}
\usepackage{lmodern}
\usepackage[a4paper,margin=2.6cm]{geometry}
\usepackage{amsmath,amssymb,amsthm}
\usepackage{xcolor}
\usepackage{tikz}
\usetikzlibrary{arrows.meta,babel}
\usepackage[colorlinks=true,urlcolor=blue!50!black,linkcolor=blue!50!black]{hyperref}

\hyphenpenalty=10000 \exhyphenpenalty=10000 \emergencystretch=3em

% Colores de las figuras (los mismos que usa la página)
\definecolor{extremo}{RGB}{47,107,79}
\definecolor{azul}{RGB}{44,74,154}

% Definiciones numeradas por unidad
\theoremstyle{definition}
\newtheorem{defn}{Definición}
% Numeración por unidad: \unidad{u}{n} (después de \input{preambulo}) hace que las
% definiciones de la nota se llamen u.n, u.(n+1), ...; n es el primer número
% libre de esa unidad (ver la tabla de CONVENCIONES.md).
\newcommand{\unidad}[2]{\renewcommand{\thedefn}{#1.\arabic{defn}}\setcounter{defn}{\numexpr#2-1\relax}}

% Observaciones y erratas del video. Usar SOLO si hay algo que señalar:
%   \begin{erratas}
%     \errata{203}{3:23}{Se dice ..., pero debe decir ...}
%   \end{erratas}
% El primer argumento son los segundos (enlaza al minuto exacto del video),
% el segundo es la marca visible, el tercero es el texto.
\newenvironment{erratas}{\section*{Observaciones y erratas del video}\begin{itemize}}{\end{itemize}}
\newcommand{\errata}[3]{%
  \edef\errurl{https://www.youtube.com/watch?v=\videoid\string&t=#1}%
  \item[\textbf{\expandafter\href\expandafter{\errurl}{#2}}] #3}
) hace que las
% definiciones de la nota se llamen u.n, u.(n+1), ...; n es el primer número
% libre de esa unidad (ver la tabla de CONVENCIONES.md).
\newcommand{\unidad}[2]{\renewcommand{\thedefn}{#1.\arabic{defn}}\setcounter{defn}{\numexpr#2-1\relax}}

% Observaciones y erratas del video. Usar SOLO si hay algo que señalar:
%   \begin{erratas}
%     \errata{203}{3:23}{Se dice ..., pero debe decir ...}
%   \end{erratas}
% El primer argumento son los segundos (enlaza al minuto exacto del video),
% el segundo es la marca visible, el tercero es el texto.
\newenvironment{erratas}{\section*{Observaciones y erratas del video}\begin{itemize}}{\end{itemize}}
\newcommand{\errata}[3]{%
  \edef\errurl{https://www.youtube.com/watch?v=\videoid\string&t=#1}%
  \item[\textbf{\expandafter\href\expandafter{\errurl}{#2}}] #3}
.
% Compilar con LuaLaTeX:  lualatex nombre-de-la-nota.tex
\usepackage[spanish,provide=*]{babel}
\usepackage[T1]{fontenc}
\usepackage{lmodern}
\usepackage[a4paper,margin=2.6cm]{geometry}
\usepackage{amsmath,amssymb,amsthm}
\usepackage{xcolor}
\usepackage{tikz}
\usetikzlibrary{arrows.meta,babel}
\usepackage[colorlinks=true,urlcolor=blue!50!black,linkcolor=blue!50!black]{hyperref}

\hyphenpenalty=10000 \exhyphenpenalty=10000 \emergencystretch=3em

% Colores de las figuras (los mismos que usa la página)
\definecolor{extremo}{RGB}{47,107,79}
\definecolor{azul}{RGB}{44,74,154}

% Definiciones numeradas por unidad
\theoremstyle{definition}
\newtheorem{defn}{Definición}
% Numeración por unidad: \unidad{u}{n} (después de % Preámbulo compartido de las notas del curso de topología algebraica.
% Cada nota empieza con \documentclass[11pt]{article}, define \videoid{...}
% (el identificador del video de YouTube) y luego hace % Preámbulo compartido de las notas del curso de topología algebraica.
% Cada nota empieza con \documentclass[11pt]{article}, define \videoid{...}
% (el identificador del video de YouTube) y luego hace \input{preambulo}.
% Compilar con LuaLaTeX:  lualatex nombre-de-la-nota.tex
\usepackage[spanish,provide=*]{babel}
\usepackage[T1]{fontenc}
\usepackage{lmodern}
\usepackage[a4paper,margin=2.6cm]{geometry}
\usepackage{amsmath,amssymb,amsthm}
\usepackage{xcolor}
\usepackage{tikz}
\usetikzlibrary{arrows.meta,babel}
\usepackage[colorlinks=true,urlcolor=blue!50!black,linkcolor=blue!50!black]{hyperref}

\hyphenpenalty=10000 \exhyphenpenalty=10000 \emergencystretch=3em

% Colores de las figuras (los mismos que usa la página)
\definecolor{extremo}{RGB}{47,107,79}
\definecolor{azul}{RGB}{44,74,154}

% Definiciones numeradas por unidad
\theoremstyle{definition}
\newtheorem{defn}{Definición}
% Numeración por unidad: \unidad{u}{n} (después de \input{preambulo}) hace que las
% definiciones de la nota se llamen u.n, u.(n+1), ...; n es el primer número
% libre de esa unidad (ver la tabla de CONVENCIONES.md).
\newcommand{\unidad}[2]{\renewcommand{\thedefn}{#1.\arabic{defn}}\setcounter{defn}{\numexpr#2-1\relax}}

% Observaciones y erratas del video. Usar SOLO si hay algo que señalar:
%   \begin{erratas}
%     \errata{203}{3:23}{Se dice ..., pero debe decir ...}
%   \end{erratas}
% El primer argumento son los segundos (enlaza al minuto exacto del video),
% el segundo es la marca visible, el tercero es el texto.
\newenvironment{erratas}{\section*{Observaciones y erratas del video}\begin{itemize}}{\end{itemize}}
\newcommand{\errata}[3]{%
  \edef\errurl{https://www.youtube.com/watch?v=\videoid\string&t=#1}%
  \item[\textbf{\expandafter\href\expandafter{\errurl}{#2}}] #3}
.
% Compilar con LuaLaTeX:  lualatex nombre-de-la-nota.tex
\usepackage[spanish,provide=*]{babel}
\usepackage[T1]{fontenc}
\usepackage{lmodern}
\usepackage[a4paper,margin=2.6cm]{geometry}
\usepackage{amsmath,amssymb,amsthm}
\usepackage{xcolor}
\usepackage{tikz}
\usetikzlibrary{arrows.meta,babel}
\usepackage[colorlinks=true,urlcolor=blue!50!black,linkcolor=blue!50!black]{hyperref}

\hyphenpenalty=10000 \exhyphenpenalty=10000 \emergencystretch=3em

% Colores de las figuras (los mismos que usa la página)
\definecolor{extremo}{RGB}{47,107,79}
\definecolor{azul}{RGB}{44,74,154}

% Definiciones numeradas por unidad
\theoremstyle{definition}
\newtheorem{defn}{Definición}
% Numeración por unidad: \unidad{u}{n} (después de % Preámbulo compartido de las notas del curso de topología algebraica.
% Cada nota empieza con \documentclass[11pt]{article}, define \videoid{...}
% (el identificador del video de YouTube) y luego hace \input{preambulo}.
% Compilar con LuaLaTeX:  lualatex nombre-de-la-nota.tex
\usepackage[spanish,provide=*]{babel}
\usepackage[T1]{fontenc}
\usepackage{lmodern}
\usepackage[a4paper,margin=2.6cm]{geometry}
\usepackage{amsmath,amssymb,amsthm}
\usepackage{xcolor}
\usepackage{tikz}
\usetikzlibrary{arrows.meta,babel}
\usepackage[colorlinks=true,urlcolor=blue!50!black,linkcolor=blue!50!black]{hyperref}

\hyphenpenalty=10000 \exhyphenpenalty=10000 \emergencystretch=3em

% Colores de las figuras (los mismos que usa la página)
\definecolor{extremo}{RGB}{47,107,79}
\definecolor{azul}{RGB}{44,74,154}

% Definiciones numeradas por unidad
\theoremstyle{definition}
\newtheorem{defn}{Definición}
% Numeración por unidad: \unidad{u}{n} (después de \input{preambulo}) hace que las
% definiciones de la nota se llamen u.n, u.(n+1), ...; n es el primer número
% libre de esa unidad (ver la tabla de CONVENCIONES.md).
\newcommand{\unidad}[2]{\renewcommand{\thedefn}{#1.\arabic{defn}}\setcounter{defn}{\numexpr#2-1\relax}}

% Observaciones y erratas del video. Usar SOLO si hay algo que señalar:
%   \begin{erratas}
%     \errata{203}{3:23}{Se dice ..., pero debe decir ...}
%   \end{erratas}
% El primer argumento son los segundos (enlaza al minuto exacto del video),
% el segundo es la marca visible, el tercero es el texto.
\newenvironment{erratas}{\section*{Observaciones y erratas del video}\begin{itemize}}{\end{itemize}}
\newcommand{\errata}[3]{%
  \edef\errurl{https://www.youtube.com/watch?v=\videoid\string&t=#1}%
  \item[\textbf{\expandafter\href\expandafter{\errurl}{#2}}] #3}
) hace que las
% definiciones de la nota se llamen u.n, u.(n+1), ...; n es el primer número
% libre de esa unidad (ver la tabla de CONVENCIONES.md).
\newcommand{\unidad}[2]{\renewcommand{\thedefn}{#1.\arabic{defn}}\setcounter{defn}{\numexpr#2-1\relax}}

% Observaciones y erratas del video. Usar SOLO si hay algo que señalar:
%   \begin{erratas}
%     \errata{203}{3:23}{Se dice ..., pero debe decir ...}
%   \end{erratas}
% El primer argumento son los segundos (enlaza al minuto exacto del video),
% el segundo es la marca visible, el tercero es el texto.
\newenvironment{erratas}{\section*{Observaciones y erratas del video}\begin{itemize}}{\end{itemize}}
\newcommand{\errata}[3]{%
  \edef\errurl{https://www.youtube.com/watch?v=\videoid\string&t=#1}%
  \item[\textbf{\expandafter\href\expandafter{\errurl}{#2}}] #3}
) hace que las
% definiciones de la nota se llamen u.n, u.(n+1), ...; n es el primer número
% libre de esa unidad (ver la tabla de CONVENCIONES.md).
\newcommand{\unidad}[2]{\renewcommand{\thedefn}{#1.\arabic{defn}}\setcounter{defn}{\numexpr#2-1\relax}}

% Observaciones y erratas del video. Usar SOLO si hay algo que señalar:
%   \begin{erratas}
%     \errata{203}{3:23}{Se dice ..., pero debe decir ...}
%   \end{erratas}
% El primer argumento son los segundos (enlaza al minuto exacto del video),
% el segundo es la marca visible, el tercero es el texto.
\newenvironment{erratas}{\section*{Observaciones y erratas del video}\begin{itemize}}{\end{itemize}}
\newcommand{\errata}[3]{%
  \edef\errurl{https://www.youtube.com/watch?v=\videoid\string&t=#1}%
  \item[\textbf{\expandafter\href\expandafter{\errurl}{#2}}] #3}

\unidad{2}{1} % unidad 2; primer enunciado: 2.1

\title{Grupos libres}
\author{Luis Jorge Sánchez Saldaña\\[2pt]\small Facultad de Ciencias, UNAM}
\date{}

\begin{document}
\maketitle

\noindent\small Notas del video
\href{https://www.youtube.com/watch?v=N-8tx7lDRBY}{\guillemotleft{}Grupos libres\guillemotright{}}
(curso de Introducción a la Topología Algebraica). Se construye el grupo libre $F_S$ en un conjunto $S$ a partir de palabras en el alfabeto $S\cup\bar S$, se definen las palabras reducidas y se enuncia la propiedad universal de los grupos libres.
\normalsize

\section*{Motivación}

En álgebra lineal, todo espacio vectorial $V$ tiene una base, es decir, un conjunto de generadores $S\subseteq V$ con la siguiente propiedad: para definir una transformación lineal de $V$ en otro espacio vectorial $W$ basta decir a dónde va la base, y luego se extiende de manera lineal. Esa es la filosofía detrás de las bases.

En el caso de grupos la cosa no es tan bonita: no cualquier grupo tiene una base. Pero sí se pueden construir ciertos grupos que tienen bases en ese sentido. Son los \emph{grupos libres}, que se denotan usualmente $F_S$, donde $S$ es el conjunto de generadores: para definir un homomorfismo de $F_S$ en algún otro grupo $G$ basta decir a dónde va $S$, y el homomorfismo se extiende a todo el grupo. En este video se construyen esos grupos.

\section*{Palabras}

Se comienza con un conjunto $S$, que será el conjunto de generadores. En el video se trabaja con un conjunto finito, solo para no complicar la notación, pero toda la construcción funciona para cualquier conjunto $S$, finito o infinito. Sea entonces
\[
  S=\{a_1,a_2,\dots,a_n\}.
\]
A partir de ahora los elementos de $S$ se consideran como \emph{letras} de un abecedario, es decir, como símbolos formales; la idea es concatenarlos para formar palabras. Antes de eso, se considera otro conjunto
\[
  \bar S=\{a_1^{-1},a_2^{-1},\dots,a_n^{-1}\},
\]
formado por los \emph{inversos formales} de los elementos de $S$, también como letras o símbolos formales. Al construir el grupo libre, estas letras van a servir tal cual como los inversos en el grupo, pero por ahora son solo símbolos formales.

Una \emph{palabra} en el alfabeto $S\cup\bar S$ se obtiene de concatenar letras del alfabeto: se pueden concatenar las que se quiera, repitiendo las que se quiera, siempre y cuando sea un número finito de ellas. Algunos ejemplos:
\begin{itemize}
  \item La \emph{palabra vacía}, que no tiene ni una sola letra; se denota $\varnothing$.
  \item Las palabras con una sola letra: $a_1$, $a_2$, $a_1^{-1}$, \dots
  \item $a_1a_3a_1^{-1}a_1a_4a_4$. Note que se pueden repetir letras: $a_4$ aparece dos veces y $a_1$ también.
  \item $a_1a_1a_1a_1a_1a_1$; la notación usual para esta palabra es $a_1^{6}$.
  \item $a_2^{-1}a_2^{-1}a_2^{-1}a_2^{-1}$, que se denota $a_2^{-4}$.
\end{itemize}

Sea
\[
  P_S=\{\text{palabras en } S\cup\bar S\}
\]
el conjunto de todas las palabras; recuerde que todas son de longitud finita. En $P_S$ hay una operación binaria: la \emph{concatenación}. Si $w_1,w_2\in P_S$, su producto es la palabra $w_1w_2$, que consiste en $w_1$ seguida de $w_2$. Por ejemplo (suponiendo que el alfabeto tiene al menos cinco letras), el producto de $a_1a_2a_3$ y $a_1^{-1}a_4a_5$ es
\[
  a_1a_2a_3a_1^{-1}a_4a_5.
\]
Esta operación tiene un neutro, la palabra vacía, pues
\[
  w\varnothing=w,\qquad \varnothing w=w
\]
para toda palabra $w$, y claramente es asociativa.

Con esto, $P_S$ tiene buenas posibilidades de ser un grupo. El problema es que no hay inversos: si $w$ es una palabra no vacía, al concatenarle otra palabra se le pegan más letras y nunca se obtiene la palabra vacía. Este problema se corrige fácilmente imponiendo una relación de equivalencia en $P_S$.

\section*{El grupo libre}

Se considera la relación de equivalencia en $P_S$ generada por
\[
  w_1\,a\,a^{-1}\,w_2\sim w_1w_2\qquad\text{para toda } a\in S,
\]
donde $w_1$ y $w_2$ son palabras. De manera intuitiva, en esta relación se puede cancelar una letra con su inverso formal: después de aplicar la relación, el inverso formal se convierte en un inverso de verdad.

\begin{defn}[Grupo libre]\label{defn:grupo-libre}
El \emph{grupo libre} en el conjunto $S$ es el conjunto de clases de equivalencia
\[
  F_S=P_S/{\sim}.
\]
\end{defn}

En efecto, $F_S$ es un grupo:
\begin{itemize}
  \item La operación de $P_S$ se hereda a $F_S$: para multiplicar dos clases se multiplican (es decir, se concatenan) los representantes,
  \[
    [w_1][w_2]=[w_1w_2].
  \]
  \item Como la operación en $P_S$ es asociativa, la operación en $F_S$ también lo es: la asociatividad se hereda.
  \item El neutro es la clase de la palabra vacía, $[\varnothing]$.
  \item Ahora sí hay inversos. Si $s_1s_2\cdots s_m$ es una palabra, con cada $s_i$ una letra del alfabeto, su inverso es el que se usa siempre en teoría de grupos: se voltea la palabra y se invierte cada letra,
  \[
    [s_1s_2\cdots s_m]^{-1}=[s_m^{-1}\cdots s_2^{-1}s_1^{-1}].
  \]
  Esto es muy fácil de verificar: al multiplicar la palabra por su inverso, por la definición de la relación de equivalencia se pueden ir cancelando cada una de las letras, y al final queda la clase de la palabra vacía.
\end{itemize}

\subsection*{Ejemplos}

\begin{itemize}
  \item Si $S=\varnothing$, el alfabeto es vacío y la única palabra es la palabra vacía. Por lo tanto $F_S=\{1\}$ es el grupo con un solo elemento.
  \item Si $S=\{a\}$ consiste de una sola letra, entonces $F_S\cong\mathbb{Z}$, el grupo cíclico infinito. Esto es porque todas las palabras que se pueden formar son, salvo equivalencia, potencias de $a$:
  \[
    F_S=\{[a^m] : m\in\mathbb{Z}\}.
  \]
\end{itemize}
Con más letras ya es más difícil describir estos grupos; son un poco más misteriosos o exóticos.

\section*{Palabras reducidas}

\begin{defn}[Palabra reducida]\label{defn:reducida}
Una palabra $w$ es \emph{reducida} si no tiene subpalabras de la forma $aa^{-1}$.
\end{defn}

Por ejemplo, la palabra $a_1a_2a_2^{-1}a_3$ no es reducida, porque en medio aparece una letra junto a su inverso. Para reducirla se cancelan esas dos letras y queda $a_1a_3$, que sí es reducida.

\medskip\noindent\textbf{Ejercicio.} Toda palabra es equivalente a una única palabra reducida; es decir, todo elemento $[w]\in F_S$ tiene un único representante reducido.

\section*{La propiedad universal}

El video termina con otro ejercicio que en realidad es un teorema, y es una de las propiedades principales de los grupos libres.

\medskip\noindent\textbf{Ejercicio.} $F_S$ satisface la siguiente propiedad universal. Dados un grupo $G$ y una \emph{función} $f\colon S\to G$, existe un único homomorfismo $\bar f\colon F_S\to G$ que extiende a $f$, es decir, tal que el siguiente diagrama conmuta:
\begin{center}
\begin{tikzpicture}[x=1cm,y=1cm]
  \node (S) at (0,1.6) {$S$};
  \node (G) at (2.2,1.6) {$G$};
  \node (F) at (0,0) {$F_S$};
  \draw[-{Stealth[length=5pt]}] (S) -- node[above] {\small $f$} (G);
  \draw[{Hooks[right,length=3pt]}-{Stealth[length=5pt]}] (S) -- (F);
  \draw[dashed,-{Stealth[length=5pt]}] (F) -- node[below right] {\small $\bar f$} (G);
\end{tikzpicture}
\end{center}
Note que $f$ es solo una función: no se le pide absolutamente nada extra, ni siquiera tendría sentido pedir que sea un homomorfismo, porque $S$ es solamente un subconjunto.

Esto es lo que se comentaba al inicio: dondequiera que se mande la base, es decir, el conjunto libre de generadores $S$, se puede extender a un homomorfismo. Para definir un homomorfismo de un grupo libre, lo único que hay que hacer es decir a dónde va la base.

\begin{erratas}
  \errata{405}{6:45}{La relación de equivalencia se escribe $w_1\,a\,a^{-1}\,w_2\sim w_1w_2$ para toda $a\in S$, pero debe ser para toda $a\in S\cup\bar S$, con la convención $(a^{-1})^{-1}=a$; es decir, también se cancela $a^{-1}a$. Si solo se cancelara $aa^{-1}$, la palabra $a^{-1}a$ no sería equivalente a la palabra vacía y $[a^{-1}]$ no tendría inverso; además, la fórmula del inverso $[s_1\cdots s_m]^{-1}=[s_m^{-1}\cdots s_1^{-1}]$ usa ambas cancelaciones.}
  \errata{576}{9:36}{Se dice «si el conjunto generador consistiera de una sola palabra», pero debe decir «de una sola letra» ($S=\{a\}$, como está en el pizarrón).}
  \errata{625}{10:25}{Por la misma razón, en la definición de palabra reducida debe decir que $w$ no tiene subpalabras de la forma $aa^{-1}$ con $a\in S\cup\bar S$, es decir, ni $aa^{-1}$ ni $a^{-1}a$ con $a\in S$.}
\end{erratas}

\end{document}
